\section{Discussion}
\label{sec:discussion}

\subsection{Why predictive validity and RL training utility diverge}
\label{ssec:dis_divergence}

The argument behind Section~\ref{ssec:res1_failure_mode} is structural
rather than statistical. Predictive validation tests the surrogate along
trajectories that always start from a real \boptest{} state and forecast
ahead for a fixed horizon; residual errors are bounded by the
horizon's length and average out across many such trajectories. Closed-loop
RL training, in contrast, tests the surrogate along trajectories that
are entirely generated by the surrogate itself: each rollout step takes
the surrogate's previous output as its own next input. Residual errors
that look harmless under predictive validation accumulate into a
self-reinforcing loop under closed-loop training, because the policy
learns to compensate for them as if they were real disturbances. When
the same policy is then deployed on \boptest, those compensations
become unmodelled inputs and the closed loop drifts. This is not a
deficiency of any particular calibration choice; it is a boundary
condition on what predictive fidelity can deliver, and it is why a
soft-regularisation construction --- which uses the calibrated twin only
to penalise actions, not to generate trajectories --- recovers
deployable performance where a fidelity-first strategy does not.

\subsection{Why physical regularization is controller-family-specific}
\label{ssec:dis_controller_family}

Three mechanisms, one per family, explain the optimal anchor values
summarised in Table~\ref{tab:lambda_summary}.
\begin{enumerate}
    \item \emph{Thermostatic PPO}. The flat, low-dimensional
    observation gives the policy limited intrinsic capacity to detect
    when its supply-temperature command is physically extreme. An
    external physical anchor binds in exactly the regime where the
    policy lacks structural cues, and $\lambda_{\mathrm{temp}}=0.10$
    is the value at which the binding is informative without crowding
    out exploration.
    \item \emph{HDRL}. The high-level setpoint layer already expresses a
    trajectory-level prior about where the zone temperature should be.
    Adding an explicit temperature anchor on top supplies redundant
    supervision that the low-level layer cannot reconcile, because the
    two loss signals act on different policy heads with different
    update cadences. Hierarchical decomposition is, in this sense, an
    architectural form of physical regularisation; the explicit
    anchor only competes with it.
    \item \emph{17-D MORL}. The rich observation stack
    (forecast features, history features, smoothed delta features)
    encodes sufficient information for the policy to detect physical
    implausibility through observation cues alone. The explicit
    penalty becomes redundant and noisy; the $5$-D-to-$17$-D
    transition in Section~\ref{ssec:res2_morl} is the cleanest
    single-axis demonstration of this effect in the paper.
\end{enumerate}
The common thread is that the value of an auxiliary physical loss
depends on what the controller can already see and on how its own
architecture already constrains its policy space. A single anchor
strength cannot fit all architectures because the architectures
differ in their implicit regularisation budgets.

\subsection{Why observation geometry matters}
\label{ssec:dis_observation}

The $5$-D-to-$17$-D MORL transition is the cleanest demonstration in
the paper of the observation-geometry effect. The backend, the
hybrid loss, and the training pipeline are held identical; only the
observation interface differs, and yet the yearly composite KPI
moves from $\msmetric=1.046$ to $\msmetric=0.099$. The largest
individual contributors, according to the feature-significance
analysis of Supplementary Table~S4, are the $4$-hour weather forecast
window and the smoothed history of supply-temperature commands. Both
add information the policy cannot otherwise infer in a stochastic
weather setting: the forecast lets the policy pre-heat in
anticipation of cold spells rather than reacting to them, and the
smoothed history disambiguates the actuator state from the zone
state. The Supplementary Table~S5 independence checks confirm that
the added channels are not redundant with the existing inputs.

\subsection{Threats to validity}
\label{ssec:dis_threats}

\begin{itemize}
    \item \textbf{Single weather file.} All training and validation use the
    \boptest{} \texttt{bestest\_air} TMY weather file. Climate-zone
    generalization is not tested.
    \item \textbf{Single building testcase.} The fidelity result applies to
    \texttt{bestest\_air}. Cross-testcase generalization is addressed (in
    the strong version of the paper) by Section~\ref{sec:results_transfer}.
    \item \textbf{KPI specificity.} The composite metric $\msmetric$ is a
    \boptest{}-specific KPI; comparisons to non-\boptest{} literature are
    limited.
    \item \textbf{Seed budget.} Canonical thermostatic, MORL and pure
    \texttt{v3} results are reported at $N=3$ seeds; HDRL is reported at
    $N=1$ seed.
    \item \textbf{Hyperparameter sensitivity.} The $\lambda_{\mathrm{temp}}$
    sweep is targeted, not exhaustive (Supplementary Table~S10). We do not
    claim global optimality of the reported values.
    \item \textbf{Hybrid recipe transferability.} Without
    Section~\ref{sec:results_transfer}, the recipe is established only on a
    single testcase; transferability is an open question.
    \item \textbf{Speed-up comparison scope.} The $85\times$ throughput
    advantage reported in Section~\ref{ssec:exp_runtime} is measured
    against the standard \boptest{} RTE HTTP--Docker deployment used in
    production benchmarking, not against bare FMU evaluation. The
    component of this figure attributable to HTTP, JSON, and Docker
    overhead versus intrinsic Modelica simulation cost is not separately
    quantified here, because the \boptest{} \texttt{bestest\_air} FMU
    ships only with a Linux \texttt{x86\_64} binary and the measurement
    host is Windows. A reproducible direct-FMU benchmark script with
    WSL2 / Docker instructions is committed
    (\texttt{evaluation/build\_speed\_benchmark\_fmu\_direct.py}); a
    third Table~\ref{tab:speed} row isolating the deployment-overhead
    component is left as a deferred reproducibility item.
\end{itemize}
